% ============================================================
%  Глава 3 — Model Predictive Control (MPC)
% ============================================================
\chapter{Прогнозирующее управление (MPC)}
\label{ch:mpc}

ЛКР, рассмотренный в главе~\ref{ch:synthesis}, даёт оптимальный регулятор
\textbf{без учёта ограничений} на управление. Реальный робот SAMURAI имеет
жёсткие физические лимиты:
\begin{equation}
  |u_v| \leq u_v^{\max} = 0{,}3 \text{ м/с}, \quad
  |u_\omega| \leq u_\omega^{\max} = 2{,}0 \text{ рад/с}.
  \label{eq:robot_limits}
\end{equation}
Если ЛКР рассчитывает $u_v = 0{,}5$ м/с, актуатор насытит на $0{,}3$ ---
возникает анти-винд-ап эффект, и оптимальность теряется.

\textbf{Model Predictive Control} (MPC) решает эту задачу: на каждом
шаге решается \emph{оптимизационная задача} с явным учётом ограничений
на горизонте $N$ шагов вперёд.

\section{Идея метода}
\label{sec:mpc_idea}

В момент времени $t = kT_s$:
\begin{enumerate}
  \item Измеряется текущее состояние $\vect{x}_k$ (или его оценка
        $\hat{\vect{x}}_k$ от наблюдателя).
  \item Решается задача оптимизации: найти последовательность
        $\vect{u}_k, \vect{u}_{k+1}, \ldots, \vect{u}_{k+N-1}$,
        минимизирующую функционал на горизонте.
  \item Применяется только \emph{первый} элемент $\vect{u}_k$.
  \item На следующем шаге $k+1$ горизонт сдвигается ---
        \textbf{принцип сдвигаемого горизонта} (receding horizon).
\end{enumerate}

\begin{notebox}
MPC можно рассматривать как ЛКР с конечным горизонтом плюс ограничения.
При $N \to \infty$ и без ограничений MPC сводится к ЛКР с тем же
матрицей $\mat{K}$.
\end{notebox}

\section{Математическая постановка}
\label{sec:mpc_formulation}

Для дискретной модели \eqref{eq:discrete_model}:

\begin{formulabox}[title={Задача MPC на горизонте $N$}]
\begin{align}
  \min_{\{\vect{u}_i\}_{i=k}^{k+N-1}}
  &\quad J = \sum_{i=k}^{k+N-1}\!\!\big(\vect{x}_i\T\mat{Q}\vect{x}_i + \vect{u}_i\T\mat{R}\vect{u}_i\big)
            + \vect{x}_{k+N}\T\mat{P}_f\,\vect{x}_{k+N}
  \label{eq:mpc_cost}\\[4pt]
  \text{при условиях:}
  &\quad \vect{x}_{i+1} = \mat{A}_d\vect{x}_i + \mat{B}_d\vect{u}_i,
  \quad i = k, \ldots, k+N-1,
  \label{eq:mpc_dynamics}\\
  &\quad \vect{u}_{\min} \leq \vect{u}_i \leq \vect{u}_{\max},
  \label{eq:mpc_input_constr}\\
  &\quad \vect{x}_{\min} \leq \vect{x}_i \leq \vect{x}_{\max}.
  \label{eq:mpc_state_constr}
\end{align}
\end{formulabox}

Здесь $\mat{P}_f$ --- терминальный штраф (terminal cost), обычно
выбирается как решение уравнения Риккати $\mat{P}_\infty$ ---
это гарантирует устойчивость замкнутой системы при $N$ достаточно большом.

\section{Сведение к QP}
\label{sec:mpc_qp}

Задача \eqref{eq:mpc_cost}--\eqref{eq:mpc_state_constr} ---
\textbf{квадратичная программа} (QP) по переменной
$\vect{U} = (\vect{u}_k\T, \ldots, \vect{u}_{k+N-1}\T)\T \in \mathbb{R}^{rN}$.

\subsection{Развёртка динамики}

Подставив \eqref{eq:mpc_dynamics} итеративно, выразим все будущие состояния
через текущее $\vect{x}_k$ и $\vect{U}$:
\begin{equation}
  \begin{pmatrix}\vect{x}_{k+1}\\ \vect{x}_{k+2}\\ \vdots\\ \vect{x}_{k+N}\end{pmatrix}
  =
  \underbrace{\begin{pmatrix}\mat{A}_d\\ \mat{A}_d^2\\ \vdots\\ \mat{A}_d^N\end{pmatrix}}_{\bm{\Phi}}\,\vect{x}_k
  +
  \underbrace{\begin{pmatrix}
    \mat{B}_d & 0 & \cdots & 0\\
    \mat{A}_d\mat{B}_d & \mat{B}_d & \cdots & 0\\
    \vdots & \vdots & \ddots & \vdots\\
    \mat{A}_d^{N-1}\mat{B}_d & \mat{A}_d^{N-2}\mat{B}_d & \cdots & \mat{B}_d
  \end{pmatrix}}_{\bm{\Gamma}}
  \,\vect{U}.
  \label{eq:lifting}
\end{equation}

\subsection{QP в стандартной форме}

Подставим \eqref{eq:lifting} в \eqref{eq:mpc_cost}, получим:

\begin{formulabox}[title={Стандартная QP-задача MPC}]
\begin{equation}
  \min_{\vect{U}} \;\tfrac{1}{2}\vect{U}\T \mat{H}\,\vect{U} + \vect{f}\T(\vect{x}_k)\,\vect{U}
  \quad \text{при} \quad
  \mat{G}\vect{U} \leq \vect{g}(\vect{x}_k),
  \label{eq:qp_standard}
\end{equation}
\end{formulabox}

где
\begin{equation*}
  \mat{H} = 2\big(\bm{\Gamma}\T\bar{\mat{Q}}\bm{\Gamma} + \bar{\mat{R}}\big),
  \qquad
  \vect{f}(\vect{x}_k) = 2\bm{\Gamma}\T\bar{\mat{Q}}\bm{\Phi}\,\vect{x}_k,
\end{equation*}
$\bar{\mat{Q}} = \operatorname{blkdiag}(\mat{Q}, \ldots, \mat{Q}, \mat{P}_f)$,
$\bar{\mat{R}} = \operatorname{blkdiag}(\mat{R}, \ldots, \mat{R})$.

Матрица $\mat{H}$ постоянна (не зависит от $\vect{x}_k$) и может быть
вычислена offline. Вектор $\vect{f}$ и правая часть $\vect{g}$
ограничений пересчитываются на каждом шаге.

\section{Алгоритм решения}
\label{sec:mpc_algorithm}

Для робота SAMURAI используется простой и надёжный подход:

\subsection{Без ограничений: явное решение}

Если ограничения \eqref{eq:mpc_input_constr} неактивны
(управление лежит внутри допустимой области), QP \eqref{eq:qp_standard}
имеет аналитическое решение:
\begin{equation}
  \vect{U}^* = -\mat{H}^{-1}\vect{f}(\vect{x}_k) = -\mat{K}_{\text{MPC}}\,\vect{x}_k,
  \label{eq:mpc_unconstrained}
\end{equation}
где $\mat{K}_{\text{MPC}} \in \mathbb{R}^{rN \times n}$ ---
\emph{статический} коэффициент, вычисляемый один раз. Применяется первая
строка размера $r \times n$ (первое управление $\vect{u}_k$).

\textbf{Это ровно та же $\mat{K}$, что и у ЛКР, при $N\to\infty$ и
$\mat{P}_f = \mat{P}_\infty$.} Для $N = 10\ldots 30$ результат
близок к ЛКР, но позволяет ввести ограничения.

\subsection{С ограничениями: QP-солвер}

При активных ограничениях задача \eqref{eq:qp_standard} решается
итеративно. В нашей реализации (см. главу~\ref{ch:implementation})
используется простой \textbf{проекционный метод}: после явного решения
\eqref{eq:mpc_unconstrained} результат проецируется на бокс
$[\vect{u}_{\min}, \vect{u}_{\max}]$:
\begin{equation}
  \vect{u}_k^{\text{clip}} = \max\!\big(\vect{u}_{\min},\,\min(\vect{u}_{\max}, \vect{u}_k^*)\big).
  \label{eq:projection}
\end{equation}

\begin{notebox}
Это \textbf{не оптимальное} решение QP с ограничениями (для оптимальности
нужен active-set или interior-point), но даёт корректное и устойчивое
поведение для робота. Полноценный QP-солвер (\texttt{cvxpy}, \texttt{qpOASES},
\texttt{scipy.optimize.minimize}) можно подключить опцией
\texttt{control.mpc.solver: qp} в конфиге.
\end{notebox}

\subsection{Терминальный штраф $\mat{P}_f$}

Согласно теореме об устойчивости MPC (Mayne et al., 2000), терминальный
штраф $\mat{P}_f$, выбранный как решение уравнения Риккати замкнутой
системы $\mat{A}_d - \mat{B}_d\mat{K}_\infty$, гарантирует устойчивость
MPC при любом $N \geq 1$. Поэтому в \filepath{matlab/design\_mpc.m}
сначала решается DARE для получения $\mat{P}_f = \mat{P}_\infty$,
затем формируется QP.

\section{Параметры MPC для SAMURAI}
\label{sec:mpc_params}

\begin{center}
\begin{tabular}{lll}
\toprule
\textbf{Параметр} & \textbf{Значение} & \textbf{Комментарий}\\
\midrule
$T_s$ & $0{,}05$ с & период цикла \texttt{motor\_node} (20 Гц)\\
$N$ & $10$ & горизонт $0{,}5$ с (соизмерим с $\tau_v$)\\
$\mat{Q}$ & $\operatorname{diag}(10,10,5,1,1)$ & точность позиции/курса\\
$\mat{R}$ & $\operatorname{diag}(1,1)$ & умеренный штраф управления\\
$\mat{P}_f$ & решение DARE & гарантия устойчивости\\
$\vect{u}_{\min}$ & $(-0{,}3,\,-2{,}0)$ & физические лимиты робота\\
$\vect{u}_{\max}$ & $(+0{,}3,\,+2{,}0)$ & --- \texttt{simulator.max\_*}\\
\bottomrule
\end{tabular}
\end{center}

При $N = 10$, $n = 5$, $r = 2$ задача QP имеет $20$ переменных,
$40$ ограничений-неравенств. Время решения на Raspberry Pi 4 ---
менее 1 мс (запас порядка для 50 мс цикла).

\section{Подбор поведения и «ошибочные» режимы}
\label{sec:behavior_tuning}

Параметры MPC вынесены в \filepath{config.yaml} и могут быть изменены
без перекомпиляции. Это позволяет:

\paragraph{1. Корректные поведения}
\begin{itemize}
  \item \textbf{Точный медленный режим}: $\mat{Q}\!\uparrow$,
        $\mat{R}\!\uparrow$, $u_{\max}\!\downarrow$ ---
        робот ездит медленно, но точно.
  \item \textbf{Агрессивный режим}: $\mat{Q}_{p_x,p_y}\!\uparrow\uparrow$,
        $\mat{R}\!\downarrow$ --- быстрая отработка, рывки.
  \item \textbf{Длинный горизонт}: $N = 30$ --- видит дальше, плавнее
        огибает препятствия (но дороже по CPU).
\end{itemize}

\paragraph{2. «Ошибочные» поведения для отладки}
\begin{itemize}
  \item \textbf{Отрицательный штраф}: $\mat{Q} = -\mat{I}$ ---
        QP станет невыпуклой, регулятор будет «гнать» в нестабильность.
  \item \textbf{Завышенные лимиты}: $u_{\max} = 100$ --- ЛКР перестанет
        насыщаться, проявится колебание из-за неточностей модели.
  \item \textbf{Полюса в правой полуплоскости} (модальный режим):
        $\lambda^{\text{уст}} = +1$ --- замкнутая система неустойчива,
        робот «убегает».
  \item \textbf{Зануление обратной связи по позиции}: $\mat{Q}_{1,1}=\mat{Q}_{2,2}=0$ ---
        робот игнорирует координаты $p_x, p_y$, движется только по
        внутренним скоростям.
\end{itemize}

\begin{warnbox}
Все «ошибочные» режимы безопасны в симуляторе (\texttt{./samurai.sh sim}),
но опасны на реальном железе. Для физических испытаний:
\begin{itemize}
  \item Поднимайте робота на подставке --- колёса не должны касаться пола;
  \item Включите аппаратный аварийный стоп
        (\texttt{config.motor.collision\_guard\_stop\_m});
  \item Не превышайте $u_{\max} = 0{,}5$ м/с.
\end{itemize}
\end{warnbox}
